%% An algebraic proof of Ghosh's Lemma 3.7
%%
%% Generated by Claude Opus 5 (Anthropic) on 2026-09-16.  Published unverified in
%% AI-Generated Mathematical Explorations, exploration 001.
%%
%% This file is the AI output as generated, with two editorial changes, both made
%% by the curator and listed here in full:
%%   1. this header;
%%   2. a provenance note added after \maketitle (the \provenance block below),
%%      so that the PDF carries its own status when it travels alone; the line
%%      "% TODO: authors" of the original was removed with it.
%% The mathematical text is untouched.  The original file is reproduced verbatim
%% in ai-output.md, next to this one.

\documentclass[11pt,reqno]{amsart}

\usepackage[T1]{fontenc}
\usepackage[utf8]{inputenc}
\usepackage{amsmath,amssymb,amsthm,mathtools}
\usepackage{array,booktabs}
\usepackage[margin=1.15in]{geometry}
\usepackage[colorlinks=true,linkcolor=blue,citecolor=blue,urlcolor=blue]{hyperref}

\theoremstyle{plain}
\newtheorem{theorem}{Theorem}[section]
\newtheorem{lemma}[theorem]{Lemma}
\newtheorem{proposition}[theorem]{Proposition}
\newtheorem{corollary}[theorem]{Corollary}
\theoremstyle{definition}
\newtheorem{remark}[theorem]{Remark}

\newcommand{\PP}{\mathbb{P}}
\newcommand{\AAA}{\mathbb{A}}
\newcommand{\KK}{\mathbb{K}}
\newcommand{\bj}{\mathbf{j}}

\title[An algebraic proof of Ghosh's Lemma 3.7]
      {An algebraic proof of Ghosh's Lemma 3.7:\\
       intersection theory in place of Brouwer degree}
\author{}
\date{September 16, 2026}

\begin{document}

\begin{abstract}
In his proof of the Casas--Alvero conjecture, Ghosh \cite{Gho} establishes an
existence lemma for ``almost counterexamples'' (his Lemma 3.7) by means of
Abel--Gontcharoff polynomials and the topological Brouwer degree over
$\mathbb{C}$. This is the only analytic ingredient of the paper, and its proof
occupies pages 11--13. We observe that, in the setting where the lemma is
actually applied ---inside the downward inductive step, where the conjecture is
already assumed in degree $n+1$--- the statement follows from the projective
dimension theorem in a few lines. The resulting proof is purely algebraic, is
valid over any algebraically closed field with no hypothesis on the
characteristic, and removes every dependence on the Euclidean topology.
\end{abstract}

\maketitle

\begin{center}\small
\fbox{\begin{minipage}{0.92\textwidth}
\textbf{Provenance.} This note was generated by an AI system (Claude Opus 5,
Anthropic) on 16 September 2026, in the context of research activities of
Santiago Laplagne at the Instituto de C\'alculo, Universidad de Buenos Aires. It
is published \textbf{unverified}: no human author claims these mathematical
results, and the arguments have not been independently checked. Current status,
original AI output, and how to report an error or a verification:
\url{https://github.com/slap/ai-mathematical-explorations}
\end{minipage}}
\end{center}

\section{Notation}

We follow \cite{Gho}. Let $\KK$ be algebraically closed, $n\ge2$, and
$R_n=\KK[x_1,\dots,x_n]$. For $1\le j\le n+1$ let $\Phi_j\colon R_n\to R_n$ be
the $\KK$-algebra automorphism
\[
  \Phi_j(x_l)=x_l-x_j\ (l\neq j),\qquad \Phi_j(x_j)=-x_j,
  \qquad \Phi_{n+1}=\mathrm{id},
\]
and let $e_k$ denote the elementary symmetric polynomial of degree $k$ in
$x_1,\dots,x_n$. For $1\le i\le n$ and $1\le j\le n+1$ put
\[
  F(i,j):=\Phi_j\big(e_{n+1-i}\big),
  \qquad \deg F(i,j)=n+1-i .
\]
These forms encode, for $f(X)=X\prod_{i=1}^n(X-x_i)$ of degree $n+1$, the
condition that $f$ and its $i$-th Hasse derivative share the root $x_j$: up to
sign, $F(i,j)$ is $f_i(x_j)$ written in the roots. For a fixed choice of indices
$\bj=(j_1,\dots,j_n)$ we write $F_i:=F(i,j_i)$.

We use the reformulation \cite[Prop.~1.4]{Gho} (from \cite{Gho24}):

\begin{proposition}[\cite{Gho}]\label{prop:reform}
The Casas--Alvero conjecture holds in degree $n+1$ over $\KK$ if and only if for
every choice $\bj$ the sequence $F_1,\dots,F_n$ is a regular sequence in $R_n$.
\end{proposition}

Since the $F_i$ are $n$ forms in $n$ variables, being a regular sequence is
equivalent to being a system of parameters, that is,
\begin{equation}\label{eq:H}
  V(F_1,\dots,F_n)=\{0\}\subseteq\AAA^n
  \qquad\Longleftrightarrow\qquad
  V_+(F_1,\dots,F_n)=\varnothing\subseteq\PP^{n-1}. \tag{H}
\end{equation}

\section{The lemma}

The statement needed in \cite{Gho} (his Lemma 3.7, via the Nullstellensatz) is:
for each $l$ there exists a degree $n+1$ polynomial $f$ with $f_i(x_{j_i})=0$ for
all $i>l$ but $f_l(x_{j_l})\neq0$; equivalently,
\[
  F_l\notin\sqrt{\big(F_{l+1},\dots,F_n\big)} .
\]

\begin{lemma}\label{lem:alg}
Let $\KK$ be algebraically closed and assume \eqref{eq:H} holds for the choice
$\bj$. Then, for every $1\le l\le n$,
\[
  F_l\notin\sqrt{\big(F_{l+1},\dots,F_n\big)}\quad\text{in } R_n .
\]
\end{lemma}

\begin{proof}
Assume the contrary and set $Z:=V_+(F_{l+1},\dots,F_n)\subseteq\PP^{n-1}$.

The $F_i$ are nonzero forms (each is a linear change of variables applied to
$e_{n+1-i}\neq0$), so $Z$ is cut out by $n-l\le n-1$ hypersurfaces of
$\PP^{n-1}$. By the projective dimension theorem, $Z\neq\varnothing$ and every
irreducible component of $Z$ satisfies
\[
  \dim Z\ \ge\ (n-1)-(n-l)\ =\ l-1 .
\]

By assumption $F_l$ vanishes on $Z$, i.e.\ $Z\subseteq V_+(F_l)$, and therefore
\[
  V_+(F_l,F_{l+1},\dots,F_n)=Z,\qquad \dim V_+(F_l,\dots,F_n)\ \ge\ l-1 .
\]

Now cut with the remaining $l-1$ hypersurfaces $V_+(F_1),\dots,V_+(F_{l-1})$.
Each intersection with a hypersurface of $\PP^{n-1}$ drops the dimension by at
most $1$, and is nonempty as long as the current dimension is $\ge1$; performing
the $l-1$ cuts starting from dimension $\ge l-1$ yields
\[
  \dim V_+(F_1,\dots,F_n)\ \ge\ (l-1)-(l-1)\ =\ 0 ,
\]
so in particular $V_+(F_1,\dots,F_n)\neq\varnothing$, contradicting
\eqref{eq:H}.
\end{proof}

\begin{remark}
The proof uses nothing about the characteristic of $\KK$: only Krull's dimension
theorem and the fact that the $F_i$ are nonzero forms. In particular,
\[
  \text{CA in degree } n+1 \text{ over } \KK
  \ \Longrightarrow\
  \text{Lemma 3.7 of \cite{Gho} over } \KK,
\]
in any characteristic. The version in \cite{Gho} is unconditional but is proved
only over $\mathbb{C}$.
\end{remark}

\begin{remark}
A point of $V_+(F_1,\dots,F_n)$ is precisely a counterexample to Casas--Alvero in
degree $n+1$ whose recycled roots are prescribed by $\bj$. Informally, the proof
says that if the $l$-th datum were a consequence of the later ones, then the full
system would be overdetermined by one equation fewer, and would therefore admit a
nontrivial solution.
\end{remark}

\section{Why the conditional version suffices}

Lemma 3.7 of \cite{Gho} is used only in the proof of his Theorem 3.6 (that
$\mu_{R_n}(\mathcal I_n(\bj))=n$), and the latter only through Corollary 3.9. In
turn, Corollary 3.9 is invoked twice, in the proofs of Proposition 4.3 and of
Lemma 4.5, both inside Section 4.2, that is, under the inductive hypothesis
4.1.1: \emph{the conjecture holds in degree $n+1$}. That hypothesis is exactly
\eqref{eq:H}.

Consequently, Lemma \ref{lem:alg} can replace Lemma 3.7 without altering the
structure of the argument: it suffices to state Theorem 3.6 in its conditional
form (``under \eqref{eq:H}, $\mu_{R_n}(\mathcal I_n(\bj))=n$''). Note also that
the form required in \cite[(3.9)]{Gho} is non-membership in the \emph{ideal}
$(F_i: i>l)$, which is weaker than the non-membership in the radical proved
here.

\begin{remark}
We would suggest that the author double-check this usage point: the reading above
comes from tracing the cross-references of the preprint (v2, March 2026), not
from a full rewriting of the paper.
\end{remark}

\section{What is removed and what is gained}

\begin{itemize}
\item The Abel--Gontcharoff polynomials, the incidence variety $\mathcal V$, the
  space $\PP(\mathcal L)$ of linear surjections, the paths $\gamma$, the balls
  $B_R(\mathbf 0)$ and spheres $S_{\epsilon_i}(z_i)$, the homotopy invariance of
  the Brouwer degree and the positivity of the local holomorphic index
  (\cite[Part I, §5.4]{AGZV}) all disappear. These are essentially pages 11--13
  of \cite{Gho}.
\item The only analytic dependence disappears: the argument becomes valid over
  any algebraically closed field, not just over $\mathbb{C}$, and in any
  characteristic.
\item Remarks 3.8(1) and 3.8(2) of \cite{Gho}, which handle special cases by
  explicit witnesses and by Rolle's theorem, become subsumed.
\end{itemize}

The moral is the usual one in intersection theory: the positivity of the local
index of a holomorphic map, which \cite{Gho} uses to guarantee that the degree
does not vanish, is unnecessary if one argues with dimensions in projective
space, where the nonemptiness of an intersection of few hypersurfaces is
automatic.

\section{Computational verification}

Using SageMath 10.7, in the root-space formulation, we verified for all choices
$\bj$ and all $l$:

\begin{center}
\renewcommand{\arraystretch}{1.2}
\begin{tabular}{>{\raggedright\arraybackslash}p{0.42\textwidth}%
                >{\raggedright\arraybackslash}p{0.22\textwidth}%
                >{\raggedright\arraybackslash}p{0.22\textwidth}}
\toprule
 & $n=3$ ($d=4$) & $n=4$ ($d=5$)\\
\midrule
hypothesis \eqref{eq:H}: $\dim V(F_1,\dots,F_n)=0$ & $64/64$ branches & $625/625$ branches\\
conclusion: $F_l\notin\sqrt{(F_i:i>l)}$ & $84/84$ cases & $780/780$ cases\\
dimension bound $\dim V(F_{l+1},\dots,F_n)\ge l$ & sharp ($=l$) & sharp ($=l$)\\
cases covered by \cite[Rmk.~3.8(1)]{Gho} & $43\%$ & $23\%$\\
\bottomrule
\end{tabular}
\end{center}

The observed dimensions are exactly $l$, the minimum allowed by the dimension
theorem: the argument rests on the tight case. The last row also shows that the
proportion of cases that \cite{Gho} settles by explicit witnesses decreases with
the degree, so the step replaced here is the one carrying almost all of the
work.

\begin{thebibliography}{9}

\bibitem{AGZV} V.~I.~Arnold, S.~M.~Gusein-Zade, A.~N.~Varchenko,
  \emph{Singularities of Differentiable Maps, Volume 1}, Birkhäuser, 2012.

\bibitem{Gho24} S.~Ghosh, \emph{A finiteness result towards the Casas--Alvero
  conjecture}, arXiv:2402.18717.

\bibitem{Gho} S.~Ghosh, \emph{Proof of the Casas--Alvero conjecture},
  arXiv:2501.09272 (v2, March 2026).

\end{thebibliography}

\end{document}
